\documentclass[aspectratio=169]{beamer}
\usepackage[utf8]{inputenc}
\usepackage[T1]{fontenc}
\usepackage[finnish]{babel}
\usepackage{minted}
\usepackage{fontawesome5}
\usepackage{booktabs}
\usepackage{textcomp}
\usepackage{tcolorbox}
\tcbuselibrary{skins}
\newcommand{\infobubble}[1]{%
  \begin{tcolorbox}[
    enhanced, size=small, fontupper=\scriptsize,
    colback=blue!8, colframe=blue!40,
    arc=4pt, boxrule=0.6pt,
    left=4pt, right=4pt, top=2pt, bottom=2pt,
  ]#1\end{tcolorbox}}

\usetheme{Madrid}
\usecolortheme{seahorse}
\setbeamertemplate{navigation symbols}{}
\setbeamertemplate{footline}[frame number]

\title{Agenttiavusteinen ohjelmointi käytännössä}
\subtitle{Tutkatuotannon uudistushanke}
\author{PAK / Tutkatuotanto}
\institute{Ilmatieteen laitos}
\date{Mapa 4.5.2026}

% ─────────────────────────────────────────────────────────────────────────────
\begin{document}

\begin{frame}
  \titlepage
\end{frame}

% ─────────────────────────────────────────────────────────────────────────────
\begin{frame}{Tehtävä}
  \begin{columns}
    \begin{column}{0.55\textwidth}
      \textbf{Lähtötilanne}
      \begin{itemize}
        \item Tuotannossa C-ohjelmat + tcsh-skriptit
        \item Tuottaa \textbf{POH}- ja \textbf{LHI}-kenttiä IRIS-tutkakomposiiteista
        \item 20-vuotias täysin palvellut koodi
      \end{itemize}
      \vspace{0.5em}
      \textbf{Tavoite}
      \begin{itemize}
        \item Modernisoi Python-paketiksi
        \item Sama toiminnallisuus, luettava koodi
        \item Airflow-yhteensopiva Docker-kontti
        \item Ulostulo: ODIM HDF5 ja Cloud-Optimized GeoTIFF
      \end{itemize}
    \end{column}
    \begin{column}{0.42\textwidth}
      %\includegraphics[width=\textwidth,height=0.65\textheight,keepaspectratio]{%
      %  pipeline_diagram_placeholder}%
      % Korvaa tämä kaavakuvalla jos haluat
      \begin{block}{Datavirta}
        \small
        IRIS TOPS (45/50 dBZ)\\
        + NWP-isotermit (0\textdegree C, -20\textdegree C)\\
        $\downarrow$ \\
        POH $\in [0,1]$ \quad LHI [m]
      \end{block}
      \vspace{0.3em}
      \infobubble{%
        \textbf{TOPS} kaikuhuipun korkeus [m]\\
        \textbf{POH} Probability of Hail\\
        \textbf{LHI} Large Hail Index
      }
    \end{column}
  \end{columns}
\end{frame}

% ─────────────────────────────────────────────────────────────────────────────
\begin{frame}{Työtapa: mob programming + Claude}
  \begin{columns}
    \begin{column}{0.48\textwidth}
      \textbf{Istuntorakenne}
      \begin{itemize}
        \item 2 päivää, reilu 6 tuntia yhteensä
        \item 6 hengen ryhmä
        \item \textbf{6 minuutin kierrot} — yksi "kuljettaja" kerrallaan
        \item Koko ryhmä seuraa ja kommentoi reaaliajassa
      \end{itemize}
    \end{column}
    \begin{column}{0.48\textwidth}
      \textbf{Valmistelu ennen istuntoja}
      \begin{itemize}
        \item Vanha koodi Git-repoon
        \item Esimerkkidata valmiina
      \end{itemize}
      \vspace{0.5em}
      \textbf{Ensimmäiset tehtävät istunnossa}
      \begin{itemize}
        \item \texttt{CLAUDE.md}:n ensimmäinen versio ennen ensimmäistäkään koodiriviä
        \item Pyydettiin sopiviin askeliin jaettu toteutussuunnitelma
      \end{itemize}
    \end{column}
  \end{columns}
\end{frame}

% ─────────────────────────────────────────────────────────────────────────────
\begin{frame}{\texttt{CLAUDE.md} — agentin kontekstidokumentti}
  \begin{columns}
    \begin{column}{0.52\textwidth}
      \textbf{Mitä se sisältää}
      \begin{itemize}
        \item Projektin tavoite ja arkkitehtuuri
        \item Käytetyt kirjastot (xarray, wradlib, \ldots)
        \item Koodaustyyli (Black, type hints, \ldots)
        \item Testausvaatimukset
        \item Ilmatieteen alan käsitteet ja termit
      \end{itemize}
      \vspace{0.5em}
      \begin{exampleblock}{}
        \small
        \faLightbulb\ Agentti lukee \texttt{CLAUDE.md}:n kontekstiksi jokaisella kehotteella
      \end{exampleblock}
    \end{column}
    \begin{column}{0.45\textwidth}
      \begin{block}{Osa kuuden tunnin tuotoksia}
        \small
        \texttt{CLAUDE.md} on osa repositoriota ja
        kehittyi projektin mukana — ei kertaluonteinen valmistelu
        vaan elävä dokumentti.\\[0.5em]
      \end{block}
    \end{column}
  \end{columns}
\end{frame}

% ─────────────────────────────────────────────────────────────────────────────
\begin{frame}{Mitä 6 tunnissa syntyi}
  \begin{columns}
    \begin{column}{0.52\textwidth}
      \textbf{Toteutetut komponentit}
      \begin{enumerate}
        \item \texttt{projection/grid.py} — FIN1000-koordinaatisto
        \item \texttt{io/iris.py} — IRIS TOPS -lukija (wradlib)
        \item \texttt{io/nwp.py} — NWP-tekstilukija + interpolointi
        \item \texttt{algorithms/poh.py} — POH-algoritmi
        \item \texttt{algorithms/lhi.py} — LHI- ja THI-algoritmit
        \item \texttt{pipeline.py} — Kokonaisputkilinja
        \item \texttt{output/odim.py} — ODIM HDF5 -kirjoitus
        \item \texttt{output/geotiff.py} — Cloud-Optimized GeoTIFF
        \item Debuggausskriptit
      \end{enumerate}
    \end{column}
    \begin{column}{0.45\textwidth}
      \begin{block}{Testikattavuus}
        \centering
        \Large \textbf{158 testiä}\\[0.3em]
        \normalsize
        yksikkö- ja integraatiotestit\\
        oikeilla IRIS-tiedostoilla
      \end{block}
      \vspace{0.5em}
      \begin{block}{Koodirivejä}
        \centering
        \small
        \begin{tabular}{lr}
          \toprule
          Komponentti & Rivejä \\
          \midrule
          Lähdekoodi   & $\sim$600 \\
          Testit        & $\sim$700 \\
          \bottomrule
        \end{tabular}
      \end{block}
    \end{column}
  \end{columns}
\end{frame}

% ─────────────────────────────────────────────────────────────────────────────
\begin{frame}{Miltä yhteistyö Clauden kanssa tuntui?}
  \begin{columns}
    \begin{column}{0.48\textwidth}
      \begin{exampleblock}{\faThumbsUp\ Toimi hyvin}
        \begin{itemize}
          \item Kevyehkön speksauksen jälkeen päästiin nopeasti toteutukseen
          \item Pysyi speksatussa tyylissä ja rakenteessa
          \item Algoritmilogiikan käänteismallinnus legacy-koodista
          \item Itsenäinen virheiden korjaus
          \item Ymmärsi alan terminologian CLAUDE.md:n pohjalta
          \item Tuotteliasta ja mielekästä
        \end{itemize}
      \end{exampleblock}
    \end{column}
    \begin{column}{0.48\textwidth}
      \begin{alertblock}{\faExclamationTriangle\ Vaati huomiota}
        \begin{itemize}
          \item Vuokaavio ei koskaan valmistunut, söi paljon tokeneita
          \item Peruslisenssillä tokenit loppui kesken
          \item Pari kriittistä virhettä piti korjata käsin
          \item Erikoiset yksikkö- ja skaalausvirheet vähenivät substanssikontekstia tarkentamalla
          \item Tasapaino jouhevan etenemisen ja suuren koodimäärän syynäämisen välillä
          \item Tasapaino määrätietoisen ohjauksen ja mikromanageroinnin välttämisen välillä
        \end{itemize}
      \end{alertblock}
    \end{column}
  \end{columns}
\end{frame}

% ─────────────────────────────────────────────────────────────────────────────
\begin{frame}{Opit}
  \begin{enumerate}
    \setlength{\itemsep}{0.7em}
    \item \textbf{CLAUDE.md on tärkein yksittäinen investointi} \\
          Konteksti pysyy yhtenäisenä joka syklillä
    \item \textbf{Mob-formaatti sopii hyvin agenttityöskentelyyn} \\
          Monta silmäparia havaitsee virheet ennen niiden juurtumista
    \item \textbf{Tehokas tapa uusien välineiden ja työskentelytapojen omaksumiseen} \\
          Mob-työskentely madaltaa ratkaisevasti kynnystä mukavuusalueelta poistumiseen
    \item \textbf{Testit samaan tahtiin toteutuksen kanssa} \\
          Agentti kirjoittaa ja ajaa testit, sekä korjaa enimmät virheet itsenäisesti
    \item \textbf{Haasta ja kysy} \\
          Claude ehdottaa, ihminen päättää — aktiivinen ohjaus on välttämätöntä
    \item \textbf{Substanssiosaaminen ei häviä mihinkään} \\
          Säätutka-asiantuntemus oli kriittistä joka askeleella
  \end{enumerate}
\end{frame}

% ─────────────────────────────────────────────────────────────────────────────
\begin{frame}{Yhteenveto}
  \begin{columns}
    \begin{column}{0.55\textwidth}
      \begin{block}{Tulos}
        Toimiva Python-toteutus syntyi vanhan tuotantokoodin pohjalta \textbf{kahdessa päivässä} kuuden hengen ryhmässä.\\[0.5em]
        Ilman agenttiapua sama työ olisi kestänyt \emph{arviolta viikkoja}.
      \end{block}
      \vspace{0.7em}
      \textbf{Seuraavat askeleet}
      \begin{itemize}
        \item Containerfile ja Airflow-integraatio
        \item Validointi oikeaa tuotantodataa vastaan
        \item Projekti jatkuu mahdollisesti samalla tyylillä
      \end{itemize}
    \end{column}
    \begin{column}{0.42\textwidth}
      \begin{block}{}
        \centering
        \vspace{0.5em}
        \Large\faCode\\[0.3em]
        \normalsize
        \texttt{radar-hail-probability}\\
        \small (hailathon)\\[1em]
        \large 6h+ \faArrowRight\ kokonainen toteutus\\
        \large + 158 testiä
        \vspace{0.5em}
      \end{block}
    \end{column}
  \end{columns}
  \vspace{1em}
  \centering
  \large\textit{Kysymyksiä?}
\end{frame}

\end{document}
